%\documentclass{gtpart}
\documentclass{amsart}
\usepackage[utf8]{inputenc} % allow utf-8 input
\usepackage[T1]{fontenc}    % use 8-bit T1 fonts
\usepackage[backref=page]{hyperref}       % hyperlinks
%\usepackage{url}            % simple URL typesetting
%\usepackage{booktabs}       % professional-quality tables
\usepackage{amsfonts}       % blackboard math symbols
\usepackage{nicefrac}       % compact symbols for 1/2, etc.
\usepackage{microtype}      % microtypography

%%% Adding ORCID 

%\usepackage{academicons}
%\definecolor{orcidlogocol}{HTML}{A6CE39}
%\newcommand{\orcid}[1]{\href{https://orcid.org/#1}{\textcolor[HTML]{A6CE39}{\aiOrcid}}}

\usepackage[pagewise]{lineno}%\linenumbers
%\linenumbers

\hypersetup{
  colorlinks   = true, %Colours links instead of ugly boxes
  urlcolor     = blue, %Colour for external hyperlinks
  linkcolor    = teal, %Colour of internal links
  citecolor   = teal %Colour of citations
}

\usepackage[usenames,dvipsnames]{xcolor}
%        \newcommand{\red}[1]{{\color{red}{#1}}}
%        \newcommand{\blue}[1]{{\color{blue}{#1}}}
        \newcommand{\green}[1]{{\color{green}{#1}}}

        \newcommand{\wild}[1]{{\color{WildStrawberry}{#1}}}

%\newcommand{\rp}[1]{{\todo{\textcolor{violet}{rp: #1}}}}

\usepackage[color=teal!70!white, bordercolor=teal, textsize=footnotesize,obeyFinal]{todonotes}

\newcommand{\ps}[1]{{\todo{\textcolor{white}{ps: #1}}}}

\newcommand{\done}{\text{\rlap{$\square$}\green{$\checkmark$}}}

%\usepackage[utf8]{inputenc}
%\usepackage{amscd}

%\usepackage{fullpage}

\usepackage{marvosym}
\usepackage{amsmath}
\usepackage{amsgen}
\usepackage{amsbsy}
\usepackage{amsopn}
\usepackage{amsthm}
\usepackage{amscd}
\usepackage{amsfonts}
\usepackage{amssymb}
\usepackage{mathrsfs}

\usepackage{ifsym}

\newcommand{\iprod}{\mathbin{\lrcorner}}

%\usepackage[marginratio=1:1,height=8.5in,width=6in,tmargin=1in]{geometry}
%\usepackage[normalem]{ulem}
%\usepackage{enumerate,enumitem}
\usepackage{tikz-cd}
\usepackage{epigraph}

% TikZ manual https://mirrors.mit.edu/CTAN/graphics/pgf/base/doc/pgfmanual.pdf

\usepackage[usestackEOL]{stackengine}
\renewcommand{\epigraphsize}{\small}
\newcommand{\mytextformat}{\epigraphsize\itshape}
\newcommand{\mysourceformat}{\epigraphsize\scshape}

\renewcommand{\textflush}{flushleft} \renewcommand{\sourceflush}{flushleft}

\let\originalepigraph\epigraph 
\renewcommand\epigraph[2]{%
  \setbox0=\hbox{\stackon{\textit{\mytextformat\Longstack{#1}}}%
    {\mysourceformat\scshape\Longstack{#2}}}%
  \ifdim\wd0>.8\linewidth\wd0=.8\linewidth\fi%
  \setlength{\epigraphwidth}{\wd0}%
  \originalepigraph{\textit{#1}}{\textsc{#2}}%
}

%\usepackage[pdftex]{hyperref}
%\usepackage[usenames]{color}

\usepackage{color}

%\usepackage[usenames,dvipsnames]{xcolor}
    %    \newcommand{\red}[1]{{\color{red}{#1}}}
        \newcommand{\blue}[1]{{\color{blue}{#1}}}
        % \newcommand{\green}[1]{{\color{green}{#1}}}

\usepackage{graphicx}
%\usepackage{mathbbol}
\usepackage{comment}
%\pdfoutput=1

\usepackage{latexsym}
%\usepackage{MnSymbol}
%\DeclareFontFamily{U}{MnSymbolC}{}
%\DeclareSymbolFont{MnSyC}{U}{MnSymbolC}{m}{n}
%\DeclareFontShape{U}{MnSymbolC}{m}{n}{
%    <-6>  MnSymbolC5
%   <6-7>  MnSymbolC6
%   <7-8>  MnSymbolC7
%   <8-9>  MnSymbolC8
%   <9-10> MnSymbolC9
%  <10-12> MnSymbolC10
%  <12->   MnSymbolC12}{}
%\DeclareMathSymbol{\intprod}{\mathbin}{MnSyC}{'270}
%\DeclareMathSymbol{\intprod1}{\mathbin}{MnSyC}{'271}
%\DeclareMathSymbol{\intprod2}{\mathbin}{MnSyC}{'269}

\newcommand{\marg}[1]
  {\mbox{}\marginpar{\tiny\hspace{0pt}#1}}

\newcommand{\mc}{\mathcal}
\newcommand{\mbf}{\mathbf}
\newcommand{\mb}{\mathbb}
\newcommand{\bb}{\mathbb}
\newcommand{\mf}{\mathfrak}
\newcommand{\msf}{\mathsf}
\newcommand{\scr}{\mathscr}
\newcommand{\ds}{\displaystyle}


\newtheorem{Theorem}{Theorem}
\newtheorem{Lemma}[Theorem]{Lemma}
\newtheorem{lemma}[Theorem]{Lemma}
\newtheorem{Corollary}[Theorem]{Corollary}

\newtheorem{Proposition}[Theorem]{Proposition}
\newtheorem{Conjecture}[Theorem]{Conjecture}
\newtheorem{Definition}{Definition}
\newtheorem{Example}[Theorem]{Example}
\newtheorem{Remark}{Remark}
\newtheorem{Question}{Question}

\newcommand{\Vol}{\operatorname{Vol}}
\DeclareMathOperator{\vol}{vol}
\DeclareMathOperator{\m}{\mf{m}}
\newcommand{\pa}{\partial}
\newcommand{\del}{\partial}
\newcommand{\op}{\operatorname}
\newcommand{\abs}[1]{\left| #1 \right|}
\newcommand{\norm}[1]{\left\| #1 \right\|}
\newcommand{\inner}[1]{\left\langle #1 \right\rangle}
\newcommand{\comp}{\mbox{\tiny{o}}}
\newcommand{\QED}{{\hfill$\Box$\medskip}}
\newcommand{\set}[1]{\left\{ #1 \right\} }
\newcommand{\tensor}{\otimes}
\newcommand{\grad}{\nabla}
\newcommand{\til}{\widetilde}
\newcommand{\of}{\circ}
\newcommand{\ol}{\overline}
\newcommand{\wh}{\widehat}
\newcommand{\h}{{hyp}}
\DeclareMathOperator{\Aut}{Aut}
\DeclareMathOperator{\Aff}{Aff}
\DeclareMathOperator{\Minvol}{Minvol}
\DeclareMathOperator{\Jac}{Jac}
\DeclareMathOperator{\Hess}{Hess}
\DeclareMathOperator{\supp}{supp}
\DeclareMathOperator{\NCD}{\mathsf{CD}}
\DeclareMathOperator{\RCD}{\mathsf{RCD}}
\DeclareMathOperator{\bary}{bar}
\DeclareMathOperator{\aint}{\strokedint}
\DeclareMathOperator{\argmin}{argmin}
\DeclareMathOperator{\pre}{pre}
\DeclareMathOperator{\directsum}{\bigoplus}

%\newcommand{\R}{\bf R}
%\newcommand{\Z}{{\bf Z}}
%\newcommand{\N}{{\bf N}}
\newcommand{\CC}{{\bf C}}
\newcommand{\HH}{{\bf H}}
\newcommand{\HK}{{\bf H}_{\bf{K}}}
\newcommand{\OO}{{\bf O}}

\newcommand{\eps}{\epsilon}
\newcommand{\Ga}{\Gamma}
\newcommand{\ga}{\gamma}
\newcommand{\la}{\lambda}
\newcommand{\La}{\Lambda}
\def\is{\rm{Iso}}

\newcommand{\cout}[1]{}
\def\co{\colon\thinspace}

\newcommand{\restr}{\mathbin{\raisebox{\depth}{\scalebox{1.5}{$\llcorner$}}}}

\renewcommand{\hat}{\widehat}

\definecolor{lime}{HTML}{A6CE39}
\DeclareRobustCommand{\orcidicon}{%
    \begin{tikzpicture}
    \draw[lime, fill=lime] (0,0) 
    circle [radius=0.16] 
    node[white] {{\fontfamily{qag}\selectfont \tiny ID}};
    \draw[white, fill=white] (-0.0625,0.095) 
    circle [radius=0.007];
    \end{tikzpicture}
    \hspace{-2mm}
}


\begin{document}

% MSC 37C10 Dynamics induced by flows and semiflows , 53E50  	Flows related to symplectic and contact structures, 68Q04  	Classical models of computation (Turing machines, etc.), 68Q17  	Computational difficulty of problems (lower bounds, completeness, difficulty of approximation, etc.) 

\title[Hamiltonian and Turing complete flows on smooth $4$-manifolds]{Hamiltonian and Turing complete flows on smooth $4$-manifolds}

\author[P. Su\'arez-Serrato]{P. Su\'arez-Serrato \href{https://orcid.org/0000-0002-1138-0921}{\orcidicon}\\ \today} 

\address{University of California, Santa Barbara, Isla Vista, California}

%\address{Max Planck Institute for Mathematics and Univeristy of California, Santa Barbara} 
%\address{{\it on leave from:} Instituto de Matem\'aticas, Universidad Nacional Aut\'onoma de M\'exico UNAM, Av. Universidad 3000, Coyoac\'an, 04510, CDMX, Mexico.} 

%M\'exico Tenochtitlan}

\dedicatory{Dedicated to the memory of Will Merry.}

\begin{abstract}
We present a Turing complete, volume preserving, smooth diffeomorphism on the $4$-sphere. The construction relies on the existence of generically rank-2 Poisson structures on $S^4$. A volume preserving diffeomorphism is constructed with the required properties using the symplectic foliation of the Poisson structure. The volume is then preserved because the Poisson structures are unimodular. 
Our main result is a consequence of this construction, showing that such a Turing complete, and therefore volume preserving, smooth diffeomorphism exists on every compact smooth orientable $4$-manifold. 
\end{abstract} 

\subjclass[2010]{37C10, 53E50, 68Q04, 68Q17  }
%\keywords{Turing complete diffeomorphisms, Poisson structures, smooth 4-manifolds}

\maketitle
Can liquids think?
\medskip

Such a preposterous question could only be asked by a being
whose central processing organ floats in cerebrospinal fluid, and
is sometimes modeled as a viscoelastic material. Nevertheless,
it may well be possible to get incompressible liquids to {\it compute}
for us.

The principles behind fluid mechanics date back to the works of ancient philosophers like Aristotle. True advances in our understanding of fluids started being made thanks largely to Leonhard Euler in 1757 \cite{Euler}.
Euler observed that liquid mercury would sink in water whereas solid mercury floated. 
Intuitively, knowing that density must somehow depend on pressure rather than mass alone, Euler set about deriving a formula to calculate density based solely on pressure. 
What resulted was a series of infamously difficult differential equations known today as the Euler Equations.
 With their help, Euler accurately measured the densities of numerous substances, thus laying the foundation for modern day fluid mechanics.
 
Vladimir Igorevich Arnold showed in 1966 that the Euler equations can be regarded as the geodesic flow equations on the group of volume preserving diffeomorphisms, with the right-invariant $L^{2}$ metric defined
by the kinetic energy of the fluid \cite{Arnold66}. 
In this light, the Euler equations describe an evolution in the Lie algebra of divergence-free vector fields, tracing the geodesics in the group of volume preserving diffeomorphisms.
This pioneering contribution highlights the importance of the geometry of volume preserving diffeomorphisms and dynamics for the understanding of Euler's equations, creating the field of topological fluid dynamics.

Nikola Tesla patented a fluid valve in 1920 \cite{Tesla}, and in the 1950s and 1960s the field of fluidics engineered computational machines that brought to life many  of these ideas. 
The applications of microfluidics are now proliferating into multiple scientific and technological domains.
 With these perspectives in mind, the question of how general can a fluid computational system be arrives naturally.
%

The concept of computability has fascinated mathematicians and computer scientists since its inception, leading to many significant breakthroughs in both fields. 
One of the most famous examples of a function that is not computable by standard Turing machines appears in the Halting Problem \cite{Turing36}. 
It asks whether there exists an algorithm that determines if a given Turing machine will ever halt upon running forever. 
%
Turing completeness refers to the property of certain computational models, such as Turing machines, that allow them to compute any function that can be expressed by means of a computable procedure. 
%

A notion of Turing completeness for a map $\varphi$ on a topological space $X$ was introduced by Moore in 1990 \cite{Moore}.
It requires that for any integer $k>0$, given a Turing machine $T$, an input tape $t$, and a finite string $(t_{1}^{\ast}, \ldots , t_{k}^{\ast})$ of symbols of the alphabet of $T$, there exists an explicitly constructible point $p$ in $X$ and an open set  $U\subset X$ such that the orbit of the map through $p$ intersects $U$ if, and only if, $T$ halts with an output tape whose positions $-k, \ldots , k$ correspond to the symbols $t_{1}^{\ast}, \ldots , t_{k}^{\ast}$. 

Tao proposed a program to study the relationship between universal computability and the possible existence of singular blowups for Euler's equations of fluid motion  \cite{Tao16, Tao16-2}. 
The overarching idea is to construct a solution that blows up in finite time using controlled, programmable hydrodynamics. 
Tao also gave an example of a coercive potential-well system whose flow is Turing complete \cite{Tao17}. 
 Such dynamics contain orbits for which knowing if they enter a particular set is undecidable. 
 Moreover, explicit descriptions of solutions to the Euler equations with particular dynamics provide evidence for the {\it Turing universality} of such Euler flows \cite{Tao18}.  

Here, we contribute to the understanding of how volume preserving diffeomorphisms can model universal Turing machines, with the next result:

\begin{Theorem}\label{thm:TCvol-pres-flow-S4}
There exists a Turing complete, volume preserving diffeomorphism on $S^4$.
\end{Theorem}

 We induce a volume preserving diffeomorphism using our previous joint work with Garc\'ia-Naranjo and Vera, that explicitly describes Poisson structures (of generic rank-$2$) on smooth $4$-manifolds \cite{GSV15}. 
 Our construction differs from previously known examples, though it is inspired by the existence of area-preserving Turing complete diffeomorphisms on disks, as demonstrated by Moore.
  It increases the growing list of  volume preserving systems that can perform general computation. 
 Such a result was proven for well dynamics on $T^4$ by Tao \cite[\S 4]{Tao17}, and for every closed oriented $3$-manifold by Cardona, Miranda, Peralta-Salas, and Presas \cite{CMPP21}. 
 
 
 The main result of this paper, obtained as a consequence of Theorem \ref{thm:TCvol-pres-flow-S4}, is that a compact orientable smooth $4$-manifold also admits such a Turing complete, volume preserving diffeomorphism:
 
 \begin{Corollary} \label{cor:TCvol-pres-flow-4-mfds}
A closed, oriented, smooth $4$-manifold admits a volume preserving, Turing complete diffeomorphism. 
\end{Corollary}
 


\begin{Remark}
It is worth noting that a simple topological existence result for a volume-preserving, Turing complete dynamical system on a 4-manifold can be achieved via suspension. 
For instance, using the Turing complete Euler flows on $S^3$ constructed in \cite{CMPP21}, embed a tubular neighborhood $S^3 \times (-1,1)$ into the 4-manifold and extend the 3D volume-preserving vector field by multiplying it by a normal scalar bump function.
However, such a topological extension is artificial in that it ignores the intrinsic geometric structure of the ambient manifold. 
The significance of our construction lies in demonstrating that Turing completeness can be embedded while respecting the geometric constraints of the generically rank-2, unimodular Poisson structures from singular Lefschetz and wrinkled fibrations.
\end{Remark}

We will now explain how the use of Poisson structures on smooth 4-manifolds allows us to extend results of Moore on specific smooth 4-manifolds, to all of them.
Moore’s foundational research demonstrated that a class of dynamical systems known as generalized shifts (which are equivalent to Turing machines in computational power) can be embedded into area-preserving diffeomorphisms on disks or smooth flows in ${\bf R}^3$.
Building on this, the use of Poisson structures on smooth 4-manifolds provides a geometric framework to extend these specific results to every compact, smooth, orientable 4-manifold.

The extension process relies on the following key concepts:
 
 \begin{enumerate}
 \item Embedding Moore’s diffeomorphisms into fibers:
 Moore showed that generalized shifts are conjugate to piecewise linear maps made of area-preserving components on Cantor blocks. Later research smoothed these into Turing complete, area-preserving diffeomorphisms of a 2-disk $D$.
Here, we extend this by using generically rank-2 Poisson structures that are adapted to singular Lefschetz fibrations on 4-manifolds. 
These fibrations allow the 4-manifold to be viewed locally as a product, such as $B\times S^2$, where $B$ is a region in the base space and $S^2$ is a regular fiber. 
By identifying a disk $D$ within the $S^2$ fiber, Moore’s Turing complete diffeomorphism $d$ can be embedded into the fiber's dynamics.\\


\item Defining a volume preserving diffeomorphism via Poisson geometry:
The Poisson structure $\pi_{f}$  on the 4-manifold is crucial because it allows a lower-dimensional area-preserving diffeomorphism to be defined and extended on the entire manifold.\\
{\bf Symplectic Leaf Preservation:} Moore’s area-preserving diffeomorphism acts on a disk. 
Because diffeomorphisms can be extended tangentially along the associated symplectic leaves (the fibers) of a Poisson manifold, the diffeomorphism $d$ can be extended to the entire 4-manifold while preserving the symplectic foliation of the Poisson structure.\\
{\bf  Volume Preservation:} A critical property of these specific Poisson structures is that they are unimodular. In a unimodular Poisson manifold, there exists a global volume form. Because the extended diffeomorphism preserves the symplectic leaves and their area forms, it preserves this global volume form. Consequently, the resulting Turing complete diffeomorphism on the 4-manifold is also volume-preserving.\\


 \item Generalization to All Smooth $4$-Manifolds:
We provide a method to bridge the gap from specific manifolds (like $S^4$) to all compact orientable smooth 4-manifolds.\\

{\bf  Wrinkled Fibrations:} More generally, a singular Lefschetz fibration on a 4-manifold can be transformed into a wrinkled fibration. 
These fibrations admit compatible Poisson structures where the symplectic leaves are the fibers. 
Moreover, wrinkled fibrations can be deformed in such a way to always contain an open region where the fiber is $S^2$, so Moore’s disk-based construction can be applied to any smooth $4$-manifold to produce a Turing complete, volume preserving diffeomorphism.
Through this Poisson-geometric approach, the unpredictability and undecidability Moore identified in specific low-dimensional systems are shown to be possible within the smooth, volume-preserving discrete dynamics of any orientable smooth 4-manifold.

\end{enumerate}


The following paragraphs first review the necessary aspects of Poisson geometry, specifically focusing on unimodularity. 
Next, we will recall how Poisson structures adapted to singular Lefschetz fibrations appear and how we can use them to lift diffeomorphisms to achieve our goals.
After establishing this context, we present our proofs of Theorem \ref{thm:TCvol-pres-flow-S4} and Corollary \ref{cor:TCvol-pres-flow-4-mfds}. 



\subsection*{Notions from Poisson geometry}

Write $\frak{X}^{k}$ to denote the smooth $k$-vector fields of the smooth manifold $M$, and $X$ a vector field on $M$. Let $[\![ \cdot \, , \cdot ]\!]$ be the Schouten--Nijenhuis bracket, $[\![ \cdot \, , \cdot ]\!]:  \frak{X}^{k} \times \frak{X}^{l} \to \frak{X}^{k+l - 1}$.
 A bivector field $\pi$ then gives rise to an operator $d_{\pi}: \frak{X}^{\cdot }\to \frak{X}^{\cdot + 1},$
 via the equation $d_{\pi}(X)=[\![ \pi, X]\!]$. Recall that a bivector $\pi$ such that $[\![ \pi, \pi ]\!]=0$ is called a Poisson bivector, and $\pi $ gives rise to a Poisson bracket $\{ \cdot,  \cdot \}_{\pi}$. Informally, a Poisson structure defined by such a bivector furnishes $M$ with a foliation with singularities, and the leaves are symplectic manifolds.
 To a function $h\in C^\infty(M)$ of a Poisson manifold $M$ 
we assign the Hamiltonian vector field $X_h$, defined by $X_h(\cdot )=\{\cdot , h \}_{\pi}$. 
In the simplest case, when the Poisson bracket comes directly from a symplectic form $\omega$, then the entire symplectic manifold is seen as a single leaf, and we have $\{ \cdot , f \} := \omega( \cdot , df)$. 
Thus, we recover the implicit definition of a Hamiltonian vector field $X_h$ (with respect to $\omega$), given by the defining equation $\iota_{X_{h}}\omega = -dh$, and both are equivalent.
%
We have a bundle map $\#_{\pi}:T^{*}M\to TM$ defined by: 
\[
\beta( \#_{\pi} (\alpha)) = \pi(\alpha , \beta), \quad \forall\, \alpha, \beta \in T^{*}M
\]
On the space of differential $1$-forms $\Omega^1(M)$, the tensor $\pi$ induces the following  Lie bracket: 
% 
\[
[ \alpha , \beta ]_{\pi} = L_{\#_{\pi}(\alpha)}\beta - L_{\#_{\pi}(\beta)} \alpha - d( \pi ( \alpha , \beta)),  \quad \forall\, \alpha , \beta \in  \Omega^1(M)
\]

The bundle map $\#_{\pi}$ induces a Lie algebra homomorphism between $[ \cdot , \cdot ]_{\pi}$ and the usual Lie bracket on vector fields. 


\subsection*{Unimodular Poisson structures}

This paragraph is inspired by the seminal work of Weinstein on modular fields and classes on Poisson manifolds \cite{Wmod}.
Let ${\mathcal L}$ denote the Lie derivative of the smooth manifold $M$. 
 For a Poisson bivector $\pi$, the equality $\,d_{\pi}^2=0$ holds. Therefore, the {\it Lichnerowicz--Poisson} cochain complex exists, and leads to the definition of the {\it Poisson cohomology spaces} of $(M,\pi)$ \cite{Lich}:
 \[
 H^{k}(M,\pi):=\ker d\pi / {\rm Im}\, d\pi, \quad \, k\in \{ 0, \ldots , \dim M \}
 \]
A vector field $X$ is called a {\it Poisson vector field} if ${\mathcal L}_{X}\pi = 0$. 
Following Weinstein \cite{Wmod}, we can now define the {\it modular vector field}. Let $\Omega$ be a positive volume form for $(M,\pi)$. 
Weinstein showed that, for a smooth function $f$, the operation $f\mapsto {\rm div}_{\Omega}\#_{\pi}(df)$ is a derivation, and hence, a vector field.
 It is called the modular vector field of $(M, \pi )$, with respect to the volume form $\Omega$. 
 This vector field is Poisson, we can denote it by $Y^{\Omega}: C^{\infty}(M)\to C^{\infty}(M)$ and it can also be defined by ${\mathcal L}_{X_f}\Omega = (Y^{\Omega}f)\Omega$. 
 Observe that the first cohomology class of ${\rm div}_{\Omega}\#_{\pi}(df)$ is independent of $\Omega$, because, for a real $a$, 
 \[
 {\rm div}_{a\Omega}\#_{\pi}(d\,\cdot) = {\rm div}_{\Omega}\#_{\pi}(d\,\cdot) + {\rm div}_{a\Omega}\#_{\pi}({\rm log} \,a).
 \]
 
The Poisson cohomology class $[Y^{\Omega}]$, called the {\it modular class} of $(M,\pi)$, is a well defined element of the first Poisson cohomology group $H^{1}(M,\pi)$. 
When the modular class vanishes, $[Y^{\Omega}]=0$, the Poisson manifold $(M,\pi)$ is called {\it unimodular} \cite{KS}. 
In an orientable, unimodular, Poisson manifold there exists a volume form that is invariant under all Hamiltonian vector fields \cite{Wmod}.
For example, a symplectic manifold is unimodular, and its Liouville volume form is invariant under every Hamiltonian vector field. 
Przybysz showed that a Poisson structure defined on the algebra of polynomials in $n$ variables which has $(n-2)$ functionally independent Casimirs is unimodular \cite{Przybysz}. 
%

Consider the $2$-disc $D$, with a symplectic form $\omega_{D}$, and let $(F, \omega_{F})$ be a compact symplectic manifold. 
On the product $D\times F$ let the projections on each factor be
$$p_{D}:D\times F\to D , \,\, p_{F}:D\times F\to F.$$
A symplectic form $\omega_{D\times F}$ on $D\times F$ can be defined by pulling back $\omega_{D}$ and $\omega_{F}$, with respect to $p_{D}$ and $p_{F}$, in that order. 
Define $\omega_{D\times F} := p_{D}^{\ast}\omega_{D}+p_{F}^{\ast}\omega_{F}$.


On the particular case of a product $(D\times F, \omega_{D\times F})$ of a disc $D$ with an orientable surface $F$, consider the bundle given by the projection onto the disc, now renamed $f: D\times F \to D$. On each preimage $f^{-1}(y)=F_{y}$ of a point $y$ in $D$, call $\omega(y)$ the restriction of $\omega_{D\times F}$ to $F_{y}$. Observe that $\omega(y)= \omega_{D\times F}|_{F_{y}} =  \omega_{F}$. Let $\mu$ be an area form of $D$ that agrees with its original orientation. Then, the volume forms that are invariant under leaf-preserving volume extensions like $\overline{d}$ on $(D\times F, \omega_{D\times F})$ are of the form 
%
\begin{equation}\label{eqn:inv-dens}
\omega(y)\wedge f^{*}\mu = \omega_{F} \wedge f^{*}\mu.
\end{equation}

Abouqateb and Boucetta showed that the modular class vanishes on a regular Poisson manifold (one without singularities) when its symplectic foliation admits a compatible Riemannian metric \cite{AB03}.
Further unimodularity criteria around a singular symplectic leaf were shown by Pedroza, Velasco-Barreras, and Vorobiev \cite{PVV18}.


\subsection*{A singular Lefschetz fibration on $S^4$}

We will need a convenient description of $S^4$, which we review now.
Consider the following singular Lefschetz fibration of $f:S^4\to S^2$, discovered by Auroux, Donaldson, and Katzarkov \cite{ADK05}. 
In their construction, the total space of the fibration of $S^4$ is obtained by gluing together three open pieces. 
The first, $X_{-}\simeq T^2\times D^2$, is a trivial fibration that lies over the southern hemisphere base $D_{-}\subset S^2$. 
The second, $X_{+}\simeq S^2\times D^2$, is another trivial fibration that lies over the northern hemisphere base $D_{+}\subset S^2$. 
The third, $W$, lies over a small neighborhood $E$ of the equator on the base $S^2$, such that $W=f^{-1}(E)$. 
In their terminology, the map $f$ has {\it indefinite quadratic singularities} along a 1-dimensional submanifold $\Delta \subset W\subset S^4$. 
This means that around each point of $\Delta$ there exist local co-ordinates $(y_1, y_2, y_3, t)$ on which $\Delta = \{ y_{i}=0 \}$. 
Under $f$, the set $\Delta $ maps to the equator inside $E\subset S^2$.
Moreover, for suitable co-ordinates on the base $S^2$ the map $f$ is represented on $W$ by:
 %
 \begin{equation}\label{eqn:f-on-W}
 (y_1, y_2, y_3, t) \mapsto y_1^2 - (y_2^2 + y_3^2)/2 +it
 \end{equation}
  
 The manifold $W$ is diffeomorphic to a solid torus with a small ball removed. 
 Its boundary components are $\partial_{+}W = T^{2}\times S^1$ and $\partial_{-}W = S^{2}\times S^1$, and $W$ can be seen as a cobordism between these two $3$-manifolds. 
 Over a point $p_{\pm}$ in $f(\partial_{\pm}W )$, denote the fiber $f^{-1}(p_{\pm})$ by $F_{\pm}$. Then $F_{+}\cong T^2$, and $F_{-}\cong S^2$. 
 The transition of the fibers $F_{t}$ lying over a path that joins $p_{+}$ to $p_{-}$ reveals a class of loops that vanish, as $F_{t}$ changes from $F_{+}$ to $F_{-}$. 
 The cycle that manifests this topological change is called the {\it  vanishing cycle}. 
 
 There are multiple possible choices for the diffeomorphisms that identify the corresponding boundaries,  $\partial X_{+}$ with $\partial_{+}W$, and $\partial X_{-}$ with $\partial_{-}W$. 
 To obtain $S^4$, two gluing diffeomorphisms need to be specified, one between $X_{-}$ and $W$, and another one between $X_{+}$ and $W$.
  Denote  these gluing diffeomorphisms by $g_{\pm}:X_{\pm}\to W$.
  Attach $X_{+}$ using the non-trivial element in $\pi_1 {\rm Diff}(S^2)$, and to glue $X_{-}$, twist the fibration using a closed path of diffeomorphisms of $T^2$ that corresponds to a unit translation transverse to the vanishing cycle's direction (see \cite{ADK05}). 
Thus put together, $X_{-}\cup W$ is the complement of a closed loop in $S^1\times S^3$. 
The absent loop is isotopic to $S^1\times \{ \ast \} \subset S^1 \times S^3$, projecting non-trivially to the $S^1$-factor.
 Therefore $X_{-}\cup W\simeq S^1\times B^3$, and gluing $X_{+}=D^2\times S^2$ along the boundary as described above yields $S^4$ \cite{ADK05}.



\subsection*{Diffeomorphisms and the local fibration structure}

Next we state some facts we will need about fibrations. 
Consider a locally trivial fiber bundle $p:M\to B$ between the closed smooth manifolds $M$ and $B$. Let ${\rm Vert }\subset TM$ be the {\it vertical} sub-bundle formed by the kernels of the differential of $p$:
\[
{\rm Vert }_{x} = \ker d_{x}p = T_{x}F_{p(x)}
\]
Here $F_{b}$ is the fiber of $p$ over $b$. A {\it connection} $\Gamma$ on the fibration $p:M\to B$, in the sense of Ehresmann, means a distribution $\Gamma : x \mapsto {\rm Hor}_{x}$ in $TM$, which splits the tangent bundle,
\[
T_{x}M = {\rm Hor}_{x} \oplus {\rm Vert }_{x}
\]
and obeys the following lifting property. For each point $x$ in $M$ and each base curve $\gamma:[0,1]\to B$ starting at $b_0 = p(x_0)$, there exists an integral curve $\tilde{\gamma}: [0,1]\to  M$ of $\Gamma$ that starts at $x_0$ and covers $\gamma$. A unique {\it horizontal} lift $\gamma_{x}: [0,\varepsilon]\to M$ exists because $T_{x}M$ splits into horizontal and vertical subspaces. %

Observe that the tangent bundle of $D\times F$ splits as a product, $T(D\times F)\simeq TD\oplus TF$.

Consider an area preserving diffeomorphism $d$ on $D$. We can define a diffeomorphism $\overline{d}$ on $D\times S^2$, that operates only on the first factor, as follows.
Let $(x,y)\in D\times S^2$, and set  $\overline{d}(x,y):=(d(x), y)$.
 
Recall that on a Poisson manifold, a diffeomorphism preserving the associated symplectic foliation and its area forms is volume preserving if the manifold is unimodular.


\subsection*{Poisson structures $\pi_{f}$ from singular Lefschetz fibrations $f$}

We will now briefly describe the construction a Poisson structure $\pi_{f}$ on $S^4$ for which the singularities of the fibration $f$ described above coincide with the symplectic foliation singularities of $\pi_{f}$ \cite[Theorem 1.1]{GSV15}. 
These structures are known to be {\it unimodular} \cite{BV21}, a property that will be crucial in our argument. 
% 

Consider the singular circle $\Delta\subset W\subset S^4$, described around equation \eqref{eqn:f-on-W}.  
An expression for the Poisson tensor $\pi_{f}$ not only in a neighbourhood of a point on $\Delta$, but on the entire normal bundle of $\Delta$ is available \cite{GSV15}.
First we consider the case when the normal bundle of the singular circle is orientable.  
The following expression for $\pi_{f}$  in the normal
bundle of the singular circle inside $W$  was obtained previously \cite[Equation 3.4]{GSV15},
\begin{equation} \label{eq:pi-gamma}
\pi=k(\theta,x_1,x_2,x_3)\left (x_1 \frac{\partial}{\partial x_2} \wedge \frac{\partial}{\partial x_3} + x_2 \frac{\partial}{\partial x_1} \wedge \frac{\partial}{\partial x_3} - x_3 \frac{\partial}{\partial x_1} \wedge \frac{\partial}{\partial x_2} \right ),
\end{equation}
where $k$ is a non-vanishing real function.

Let $N_{\Delta}$ be a neighborhood of $\Delta$, whose normal bundle is orientable. So $N_{\Delta}\cong S^1\times B^3$. On $N_{\Delta}\subset W$, the bivector $\pi_{f}$  is proportional to the product Poisson structure of  $S^1$ with the 
zero Poisson structure, and $B^3$ with the Lie--Poisson structure of $\frak{sl}(2,\mathbb{R})^*$ (relative to an appropriate basis identification \cite{GSV15}).

Second, when the normal bundle of a singular  circle is not orientable, we still have a coordinate description in terms of the quotient by the action of 
\begin{equation}
\label{E:defiota}
i \colon (\theta,x_1,x_2,x_3) \to (\theta+\pi, -x_1, -x_2, x_3).
\end{equation}

The involution $i$ defined on $S^1\times B^3$ by equation \eqref{E:defiota} is a  Poisson map for the bracket  defined by \eqref{eq:pi-gamma}, provided that $k\circ i=k$ \cite[Proposition 3.2]{GSV15}.
The formula for the symplectic form on the symplectic leaf $\Sigma_q$ through 
$q=(\theta,x_1,x_2,x_3)$ at the
point $q$, induced by $\pi$, is given by (cf \cite[Proposition 3.3]{GSV15}):
\begin{equation}
\label{E:def-omega-Sigma}
\omega_{\Sigma_q}=\frac{1}{k(\theta,x_1,x_2,x_3)\sqrt{x_1^2+x_2^2+x_3^2}}\,\omega_{Area}
\end{equation}
Here, $\omega_{Area}$ is the area form on $\Sigma_q$ induced by the metric $ds^2=d\theta^2+dx_1^2+dx_2^2+dx_3^2$ on $S^1\times B^3$. 

\subsection*{Area preserving diffeomorphisms from generalized shifts}

Moore showed that any generalized shift is conjugated to a piece-wise linear map, itself made up of finitely many area-preserving components, defined on Cantor blocks, restricted to the Cantor square \cite{Moore}. % 
Moore mentioned that such a piece-wise linear map could be extended to a diffeomorphism. 
Cardona, Miranda, Peralta-Salas, and Presas described one such smoothing process, thus unveiling a Turing complete, {\it area-preserving} diffeomorphism $d$ of the $2$-disc $D$ \cite[Theorem 5.2]{CMPP21}. It equals the identity on a neighborhood of the boundary of $D$. Figures (\ref{fig:1}) and (\ref{fig:2}) illustrate the components of Moore's area preserving diffeomorphism $d$, that realizes a generalized shift smoothly.
For brevity, we recommend interested readers consult more details in these two sources we use, by Moore \cite{Moore}, and by Cardona, Miranda, Peralta-Salas, and Presas \cite{CMPP21}.

\begin{figure}%[htbp]
\centering
\includegraphics[width = \textwidth]{disc-sigma-rho.pdf}
\caption{Here we illustrate the first couple of steps in a possible parametrization of Moore's diffeomorphism to realize it in an area preserving way. First, with a clockwise $\pi$ rotation $\rho$, followed by exchanging the sectors $b$ and $d$ while preserving the area of small neighborhoods of them with the map $\sigma$. These two transformations, $\rho$ and $\sigma$, all preserve area, so their composition $\sigma\circ \rho$ is also an area preserving diffeomorphism.}
\label{fig:1}
\end{figure}

\begin{figure}%[htbp]
\centering
\includegraphics[width = .7\textwidth]{disc-stretch.pdf}
\caption{Here we show the last step in an area preserving version of Moore's diffeomorphism. The teal mass on the left side sector $a$ circulates to fill the area $\bar{a}$ shown in teal on the right. Likewise, the blue mass on the left side sector $d$ circulates to end occupying the blue area $\bar{d}$ on the lower right}
\label{fig:2}
\end{figure}

We are now ready to present the proofs of our main results.
 
\begin{proof}[Proof of Theorem \ref{thm:TCvol-pres-flow-S4}]

 Consider the singular Lefschetz fibration of Auroux--Donaldson--Katzarkov $f:S^{4}\to S^2$ explained above, assuming the single critical circle $C=f(\Delta)$ lies exactly on the equator of the base $S^{2}$. 
 %
 Let $B$ be an open set on the base such that $B\subset S^2\setminus \Delta$, such that the fibration trivializes over $B$, so that $f^{-1}(B)\simeq B\times S^2$.
 Here the second factor $S^2$ is the regular fiber of $f$ over $B$.
 Let $D$ be a disk in the fiber $S^{2}\subset f^{-1}(B)\simeq B\times S^2$.
 
 After composing with an area-preserving projection, we may assume that the area-preserving, Turing complete diffeomorphism $d$ described above is acting on $D\subset S^2$. 
 For convenience, we will still call this diffeomorphism $d$.

 Observe that the diffeomorphism $d$ equals the identity on a neighborhood of the boundary of $D$, so we can extend $d$ as the identity everywhere on the complement of $D$ in $S^2$, thus defining an area-preserving diffeomorphism on $S^2$.
 
 Moreover, the area-preserving diffeomorphism $d$ on $D$ is isotopic to the identity because it originates from smoothing a generalized shift.
 Hence there exists a smooth family of area-preserving diffeomorphisms $d_t : S^2 \to S^2$ for $t \in [0,1]$ such that $d_0 = \mathrm{id}_{S^2}$, $d_1 = d$, and every $d_t$ is the identity outside $D$.

 Let $B_0 \subset B$ be a smaller open disk inside $B$. Choose a smooth bump function $\rho: B \to [0,1]$ such that $\rho(b) = 1$ for $b \in B_0$, and $\rho$ vanishes near the boundary $\partial B$.

 We extend and define a single diffeomorphism $\overline{d}$ on $B\times S^2$ by setting $\overline{d}(b, y) = (b, d_{\rho(b)}(y))$ for $(b, y) \in B\times S^2$. 
 Because $\rho(b) = 0$ near $\partial B$, the map $\overline{d}$ equals the identity near the boundary of $B \times S^2$. Therefore, it smoothly extends to all $S^{4}$ by setting it equal to the identity on $S^{4}\setminus f^{-1}(B)$, thus defining a smooth diffeomorphism on all $S^{4}$ which we will still call $\overline{d}$. 

 Notice that the diffeomorphism $\overline{d}$ in $S^4$ only moves points along the direction tangent to the symplectic leaves (the fibers of $f$). 
 Observe that, by construction, $\overline{d}$ maps each symplectic leaf of the Poisson structure $\pi_{f}$ defined by $f$ on $S^4$ to itself, preserving its respective area form.
 Recall that the Poisson structure $\pi_{f}$ is unimodular \cite{GSV15, BV21}.
 Therefore, preserving the symplectic leaves and their area forms guarantees that the global volume form associated with the unimodular Poisson structure is invariant under $\overline{d}$ on $S^{4}$.

 We will now show the existence of a constructible point and a constructible set to verify the Turing completeness property of the diffeomorphism $\overline{d}$.
 Let $p$ be the constructible point, together with $\mathcal{O}$ its associated constructible open set, guaranteed to exist by the Turing completeness of $d$ on $D\subset S^2$.
 Consider a parametrization of $B_0$ by polar coordinates, and define $c$ to be the center of $B_0$ in these coordinates. 
 Then the point $\bar{p}:= (c,p)$ is a constructible point, for which the open set $\overline{\mathcal{O}}:= B_0 \times \mathcal{O}$ is a constructible set which together imply that the diffeomorphism $\overline{d}$ is Turing complete (since $\rho(c)=1$).
 The undecidability of the orbit of $p$ under $d$ entering $\mathcal{O}$ implies, by construction, that it is undecidable if the orbit of $\bar{p}$ under $\overline{d}$ enters $\overline{\mathcal{O}}$.

 Therefore, we have identified both a constructible point $\bar{p}$ and its corresponding open set $\overline{\mathcal{O}}$, that together verify the Turing completeness of $\overline{d}$.
 %
 This concludes the proof of Theorem \ref{thm:TCvol-pres-flow-S4}, as $\overline{d}$ is a Turing complete, volume preserving diffeomorphism of $S^4$.
\end{proof}


We now present a proof of Corollary \ref{cor:TCvol-pres-flow-4-mfds}:

\begin{proof}[Proof of Corollary \ref{cor:TCvol-pres-flow-4-mfds}]
Let $M$ be a closed oriented smooth $4$-manifold.
The proof repeats the construction as in Theorem \ref{thm:TCvol-pres-flow-S4}, except using a generically rank-$2$ Poisson structure $\pi_{f}$ adapted to a slightly different kind of singular Lefschetz fibration $f$ on $M$ \cite{GSV15, ST16}. 
In this general case, Lekili showed that it is always possible to deform a singular Lefschetz fibration into a kind of fibration that is called wrinkled, and which has additional types of singularities \cite{L}. 
These wrinkled fibrations admit Poisson structures that are compatible in the sense that the symplectic leaves are the fibers, and both structures have exactly the same singularities \cite{ST16}.
As before, explicit normal forms for $\pi_{f}$ on neighbourhoods of Lefschetz singularities are available \cite{ST16}.

After performing enough wrinkled fibration moves, we may assume we have a new fibration $\bar{f}$ from $M$ to $S^2$, which has an open region where the fibre is $S^2$.
%
Choose a disc region $B$ on the base of the fibration $\bar{f}$, so that the preimage $\bar{f}^{-1}(B)$ is a product $B\times S^2$.  
Then the arguments in Theorem \ref{thm:TCvol-pres-flow-S4} carry through exactly the same.
Specifically, by the extension arguments mentioned previously, there exists a Turing complete, volume preserving diffeomorphism with respect to the Poisson structure $\pi_{\bar{f}}$ on $M$.  
\end{proof}%-----

%In the constructions in the proofs of Theorem \ref{thm:TCvol-pres-flow-S4} and Corollary \ref{cor:TCvol-pres-flow-4-mfds}, the Hamiltonian Turing complete flows vanish in large portions of the ambient manifold. 
%Observe that if a manifold admits a non-singular flow, then it has zero Euler characteristic.
%Consequently, a general result about nonsingular flows that encompasses all smooth, closed, oriented $4$-manifolds is unattainable, as singularities would be inevitable for manifolds with nonzero Euler characteristic.\\
%
%{\bf Question 1:}  Can a Hamiltonian, Turing complete, non-singular and non-vanishing flow be constructed on smooth $4$-manifolds with zero Euler characteristic?\\

%It may nevertheless be interesting to wonder if such a non-singular, and non-vanishing, volume preserving and Turing complete flow can be constructed on smooth $4$-manifolds with vanishing Euler characteristic. 

   %-----
 
% \medskip 
% 
% The flow of a solution to a system like the one here only simulates a finite number of Turing machine steps (see \cite{CMPP21}). This phenomenon is related to the {\it Space-bounded Church Turing thesis}, introduced by Braverman, Schneider, and Rojas \cite{BSR15}. Storage space complexity and memory availability are studied as constraints for Turing machines modelled on compact spaces.  Moreover, {\it universal} Turing machines cannot be robustly simulated on a compact space \cite{BGH13}. 
%
%%-----
%
% \medskip 
%  
%  It would be interesting to relate the Hamiltonian Turing complete flows presented here to a solution of Euler's equations.  
%   Our arguments rest on a background Poisson structure. 
%   We have not assumed the existence of a Riemannian metric.  
%  Hence, a model of hydrodynamics---perhaps some generalised Euler type equations---adapted to this particular situation would need to be specified first.\\
%   
%  {\bf Question 2:} Can a model of hydrodynamics analogous to Euler's equations be defined using only a Poisson structure?
% 
%%%%%
%\subsubsection*{Acknowledgments} I am grateful to Misael Avenda\~no-Camacho for reading a previous version and suggesting improvements. Thanks to Jos\'e C. Ru\'iz Pantale\'on for useful conversations and correspondance about Poisson and Hamiltonian flows.
%Thanks also to a referee who noticed the originally proposed construction was incorrect, their thoughtful comments led to improvements in this version.
%\\
%%%%


\subsubsection*{Declarations}


{\bf Conflicts of interests:}
There are no financial or non-financial conflicts of interest that are directly or indirectly related to the work submitted for publication.\\

{\bf Data availability:} There are no data associated with this work.\\

{\bf Funding:} There is no grant or project associated with this work.

\small 

\providecommand{\bysame}{\leavevmode\hbox to3em{\hrulefill}\thinspace}

\begin{thebibliography}{}
\nocite{*}

\bibitem{AB03} A. Abouqateb, M. Boucetta, {\it The modular class of a regular Poisson manifold and the Reeb class of its symplectic foliation},  C. R. Acad. Sci. Paris, Ser. I 337 (2003) 61--66.

\bibitem{Arnold66}  V.I. Arnold, {\it Sur la géométrie différentielle des groupes de Lie de dimension infinie et ses applications à l'hydrodynamique des fluides parfaits,} Ann. Inst. Fourier 16, fa.1, (1966) 319--361.

\bibitem{ADK05} D. Auroux, S.K. Donaldson, L. Katzarkov, {\it Singular Lefschetz pencils,} Geometry \& Topology 9 (2) (2005) 1043--1114.

\bibitem{BV21} P. Batakidis, R. Vera, {\it Poisson cohomology of broken Lefschetz fibrations}, Diff. Geom. App. 72 (2020) 101661.

%\bibitem{BGH13} O. Bournez, D.S. Graça, E. Hainry, {\it Computation with perturbed dynamical systems,} Journal of Computer and System Sciences 79(5) (2013)  714--724.

%\bibitem{BSR15} M. Braverman, J. Schneider, C. Rojas, {\it Space-bounded Church-Turing thesis and computational tractability of closed systems,} Phys. Rev. Lett. 115(9) (2015): 098701.

\bibitem{CMPP21} R. Cardona, E. Miranda, D. Peralta-Salas, F. Presas, {\it Constructing Turing complete Euler flows in dimension 3,} Proceedings of the National Academy of Sciences 118 (19) (2021).

\bibitem{Euler} L. Euler,  {\it Principes généraux du mouvement des fluides,}  Mémoires de l'académie des sciences de Berlin, (1757) 11: 274--315. 

\bibitem{GSV15} L.C. García-Naranjo, P. Suárez-Serrato, R. Vera, {\it Poisson structures on smooth 4–manifolds,} Lett. Math. Phys. 105(11) (2015) 1533--1550.

\bibitem{KS} Y.~Kosmann-Schwarzbach, {\it  Poisson manifolds, Lie algebroids, modular classes: a survey,} SIGMA Symmetry Integrability Geom. Methods Appl. 4 (2008), Paper 005.

\bibitem{L}  Y. Lekili, {\it  Wrinkled fibrations on near-symplectic manifolds,} Geom. Topol. 13 (2009), no. 1, 277--318.

\bibitem{Lich} A. Lichnerowicz, {\it Les variétés de Poisson et leurs algèbres de Lie associées,} Jour. Diff. Geom. 12 (1977), no. 2, 253--300.

\bibitem{Moore} C. Moore, {\it  Unpredictability and undecidability in dynamical systems,} Phys. Rev. Lett. 64, 2354(1990).

\bibitem{PVV18} A. Pedroza, E. Velasco-Barreras, Y. Vorobiev, {\it Unimodularity criteria for Poisson structures on foliated manifolds}, Lett. Math. Phys. (2018) 108:861--882.

\bibitem{Przybysz} R. Przybysz,  {\it On one class of exact Poisson structures}, J. Math. Phys. 42 (2001), no. 4, 1913--1920.

\bibitem{ST16} P. Su\'arez-Serrato, J. Torres Orozco, {\it Poisson structures on wrinkled fibrations,} Bol. Soc. Mat. Mex.  22(1) (2016) 263--280.

\bibitem{Tesla} N. Tesla, {\it Valvular conduit}, United States Patent Office, Pat. No. 1,329,559, filed Feb. 21, 1916, issued Feb. 3, 1920.

\bibitem{Tao16} T. Tao, {\it Finite time blowup for an averaged three-dimensional Navier-Stokes equation},  J. Amer. Math. Soc. 29 (2016), no. 3, 601--674.

\bibitem{Tao16-2} \bysame , {\it  Finite time blowup for Lagrangian modifications of the three-dimensional Euler equation,} Ann. PDE 2 (2016), no. 2, Art. 9.

\bibitem{Tao17} \bysame , {\it On the universality of potential well dynamics,}
Dyn. Partial Differ. Equ. 14 (2017), no. 3, 219--238.

\bibitem{Tao18} \bysame , {\it On the universality of the incompressible Euler equation on compact manifolds}, Discrete Contin. Dyn. Syst. 38 (2018), no. 3, 1553--1565.

\bibitem{Turing36}  A.M. Turing,  {\it On Computable Numbers, with an Application to the Entscheidungsproblem}
Proc. London Math. Soc. (2) 42 (1936), no. 3, 230--265.

\bibitem{Wmod} A. Weinstein, {\it The modular automorphism group of a Poisson manifold,} J. Geom. Phys. 23 (1997), no. 3-4, 379--394.

\end{thebibliography}

\end{document}





